\chapter{Kiến thức nền áp dụng}

\section{Văn phạm phi ngữ cảnh và sinh câu}

\subsection{Biểu diễn CFG}

Văn phạm phi ngữ cảnh (Context-Free Grammar, CFG) có thể biểu diễn bằng bộ $G=(N,\Sigma,P,S)$, trong đó $N$ là tập nonterminal, $\Sigma$ là tập terminal, $P$ là tập production và $S$ là ký hiệu bắt đầu. Một production có dạng $A\rightarrow\alpha$, với $A\in N$ và $\alpha$ là chuỗi terminal/nonterminal. Một câu $w$ thuộc ngôn ngữ của văn phạm nếu có thể suy dẫn $S\Rightarrow^{*}w$ \cite{jurafsky-martin}.

Trong bài toán này, nonterminal biểu diễn loại câu hỏi hoặc thành phần câu, chẳng hạn nhóm hỏi lịch học, tham chiếu tuần/chương và phần vị ngữ câu hỏi. Terminal là các cụm từ tiếng Việt cụ thể. Grammar hiện có dùng \texttt{<S>} làm ký hiệu bắt đầu; các nhánh \texttt{<Q\_*>} tổ chức các nhóm yêu cầu trong miền môn học.

\subsection{Sinh câu từ production}

Sinh câu là quá trình bắt đầu từ $S$ rồi lần lượt thay nonterminal bằng một production cho đến khi chuỗi chỉ còn terminal. Bộ sinh trong dự án duyệt breadth-first: mỗi trạng thái chứa chuỗi ký hiệu và số lần khai triển; thuật toán chọn nonterminal đầu tiên, tạo trạng thái kế tiếp cho từng production, rồi đưa các trạng thái vào hàng đợi. Khi không còn nonterminal, chuỗi được ghép thành câu.

Breadth-first giúp duyệt các nhánh theo số bước khai triển. Do CFG có thể đệ quy và tạo số lượng câu rất lớn, chương trình giới hạn độ sâu khai triển và tổng số câu đầu ra. Câu trùng được loại bỏ để mỗi dòng trong file kết quả là một chuỗi khác nhau. Những giới hạn này là điều kiện dừng của phép liệt kê, không phải đánh giá về độ tự nhiên của câu.

\section{Truy xuất thông tin và TF-IDF}

\subsection{Knowledge base, document và chunk}

Knowledge base là tập văn bản nguồn về môn học. Để tìm bằng chứng ở mức chi tiết, mỗi tệp được chia thành các đơn vị nội dung nhỏ hơn gọi là chunk. Chunking theo ranh giới có ý nghĩa (tuần, chương, mục FAQ, mục quy định) giúp hạn chế việc trộn nội dung không liên quan. Metadata như nguồn, section, loại nội dung và entity ID cho phép truy nguyên chunk được chọn.

\subsection{TF-IDF}

TF-IDF đánh trọng số một term dựa trên tần suất term trong văn bản và mức độ hiếm của term trên toàn bộ tập văn bản. Trực giác của phép đo là một từ xuất hiện nhiều trong một chunk nhưng không phổ biến trong mọi chunk có khả năng giúp phân biệt nội dung. Trong cài đặt hiện tại, \texttt{TfidfVectorizer} tạo vector cho các chunk và query; \texttt{cosine\_similarity} so sánh hướng hai vector \cite{manning-ir, sklearn-text}.

Với vector query $\vec q$ và vector chunk $\vec d$, cosine similarity là:
\[
\operatorname{cos}(\vec q,\vec d)=
\frac{\vec q\cdot\vec d}{\|\vec q\|\,\|\vec d\|}.
\]
Điểm gần 1 biểu thị hai vector gần cùng hướng; điểm gần 0 biểu thị ít term có trọng số chung. Đây là tín hiệu lexical, không phải xác suất câu trả lời đúng.

\subsection{Semantic frame và xếp hạng kết hợp}

Khi parser nhận được câu hỏi, semantic frame có thể cung cấp intent và entity ID. Những tín hiệu này giúp ưu tiên đúng loại nguồn hoặc đúng thực thể, bổ sung cho mức trùng khớp từ vựng. Với frame có intent hợp lệ, điểm combined trong cấu hình hiện tại là:
\[
\operatorname{score}(q,d)
=0.75\,\operatorname{score}_{lexical}(q,d)
+0.25\,\operatorname{score}_{semantic}(q,d).
\]
Khi không có frame hợp lệ, retriever dùng điểm TF-IDF lexical. Entity alias trong \texttt{entities.txt} được dùng để chuẩn hóa cách diễn đạt và mở rộng query. Điểm combined và điều kiện ngưỡng được dùng để chọn tối đa $k$ chunk; nếu không có chunk phù hợp thì trả \texttt{NO\_RELEVANT\_CONTEXT}.

\section{Vai trò của retrieval trong pipeline}

Retrieval xếp hạng và cung cấp đoạn văn bản làm bằng chứng. Nó không đồng nhất với answer generation: có một chunk liên quan không đảm bảo chunk chứa thông tin cần thiết để trả lời. Đặc biệt, trường hợp query về ngày hạn nộp nhưng dữ liệu chỉ chứa yêu cầu BTL có thể truy xuất đúng phần BTL mà vẫn thiếu ngày cụ thể. Tầng trả lời phía sau cần phân biệt \emph{có context liên quan} với \emph{context đủ để trả lời}.

Các phần bổ sung trong pipeline gồm Earley parsing, clause validation cho multi-intent, semantic frame, RAG, Answer Guard và dialogue state. Earley parsing dùng chart states với Predictor, Scanner và Completer. Multi-intent chỉ được chấp nhận khi từng clause parse độc lập thành công; semantic frame giữ riêng intent/entity cho từng clause. RAG cung cấp context từ knowledge base cho LLM; Answer Guard kiểm tra context, citation và một số facts có cấu trúc. Dialogue manager giải tham chiếu bằng các pattern đã định nghĩa, không phải coreference resolution tổng quát.
